% Setting up the document with a comprehensive preamble
\documentclass[a4paper,12pt]{article}
\usepackage[utf8]{inputenc}
\usepackage[T1]{fontenc}
\usepackage{geometry}
\geometry{margin=1in}
\usepackage{amsmath, amssymb}
\usepackage{graphicx}
\usepackage{caption}
\usepackage{listings}
\usepackage{xcolor}
\usepackage{tikz}
\usepackage{tikz-uml}
\usepackage{hyperref}
\hypersetup{colorlinks=true, linkcolor=blue, urlcolor=blue}
\usepackage{fancyhdr}
\pagestyle{fancy}
\fancyhead[L]{Arbitrary-Precision Arithmetic Library}
\fancyhead[R]{\today}
\usepackage{enumitem}
\usepackage{titling}
\usepackage{noto}

% Defining colors for code listings
\definecolor{codegreen}{rgb}{0,0.6,0}
\definecolor{codegray}{rgb}{0.5,0.5,0.5}
\definecolor{codepurple}{rgb}{0.58,0,0.82}
\definecolor{backcolour}{rgb}{0.95,0.95,0.92}

% Setting up the listings package for code
\lstdefinestyle{mystyle}{
    backgroundcolor=\color{backcolour},   
    commentstyle=\color{codegreen},
    keywordstyle=\color{magenta},
    numberstyle=\tiny\color{codegray},
    stringstyle=\color{codepurple},
    basicstyle=\ttfamily\small,
    breakatwhitespace=false,         
    breaklines=true,                 
    captionpos=b,                    
    keepspaces=true,                 
    numbers=left,                    
    numbersep=5pt,                  
    showspaces=false,                
    showstringspaces=false,
    showtabs=false,                  
    tabsize=4
}
\lstset{style=mystyle}

% Title and author
\title{Arbitrary-Precision Arithmetic Library: Project Report}
\author{
   Gagan Chandra \\
   Cs24BTech11020
}
\date{May 2025}

\begin{document}

% Generating the title page
\maketitle
\tableofcontents
\newpage

% Introducing the design section
\section{Design Decisions}
Our project implements an arbitrary-precision arithmetic library in Java, supporting both integers and floating-point numbers. Below, we outline our design decisions and provide UML diagrams for the core classes.

\subsection{Design Overview}
We designed the library with the following goals:
\begin{itemize}
    \item \textbf{Arbitrary Precision}: Handle numbers of any size without loss of precision.
    \item \textbf{Simplicity}: Use a beginner-friendly approach for implementation.
    \item \textbf{Reusability}: Separate integer and float operations into distinct classes.
    \item \textbf{Automation}: Include scripts for building and running the application.
\end{itemize}

\subsection{Class Structure}
The library is organized into a package \texttt{arbitraryarithmetic} with two main classes:
\begin{itemize}
    \item \texttt{AInteger}: Handles arbitrary-precision integers.
    \item \texttt{AFloat}: Handles arbitrary-precision floating-point numbers with up to 30 decimal digits.
\end{itemize}
A separate \texttt{MyInfArith} class serves as the main application, taking command-line arguments to perform operations. We also included a Python script (\texttt{compile\_and\_run.py}) for automation and an Ant build script (\texttt{build.xml}) for compilation and packaging.

\subsection{Data Representation}
\begin{itemize}
    \item \texttt{AInteger}: Represents numbers as strings (e.g., "123456789") with a boolean \texttt{isNegative} to track the sign. This avoids Java's integer size limits.
    \item \texttt{AFloat}: Splits numbers into \texttt{integerPart} and \texttt{fractionalPart} (e.g., "123.456" as "123" and "456"), with \texttt{isNegative}. The fractional part is truncated to 30 digits to match the specification.
\end{itemize}

\subsection{Arithmetic Operations}
\begin{itemize}
    \item \texttt{AInteger}: Implements addition, subtraction, multiplication, and division using string-based algorithms (e.g., digit-by-digit addition with carry).
    \item \texttt{AFloat}: Converts floats to integers by aligning decimal points, uses \texttt{AInteger} for operations, then adjusts the decimal point in the result.
\end{itemize}

\subsection{UML Diagrams}
Below are UML class diagrams for \texttt{AInteger}, \texttt{AFloat}, and \texttt{MyInfArith}.

% Creating UML diagram for AInteger
\begin{figure}[h]
    \centering
    \begin{tikzpicture}
        \begin{umlpackage}{arbitraryarithmetic}
            \umlclass{AInteger}{
                - number: String \\
                - isNegative: boolean \\
                \\
                + AInteger() \\
                + AInteger(s: String) \\
                + AInteger(other: AInteger) \\
                + parse(s: String): AInteger \\
                + toString(): String \\
                + add(b: AInteger): AInteger \\
                + subtract(b: AInteger): AInteger \\
                + multiply(b: AInteger): AInteger \\
                + divide(b: AInteger): AInteger \\
                - removeLeadingZeros(s: String): String \\
                - compareAbsolute(a: String, b: String): int \\
                - addStrings(a: String, b: String): String \\
                - subtractStrings(a: String, b: String): String \\
                - multiplyStrings(a: String, b: String): String \\
                - divideStrings(a: String, b: String): String
            }{}
        \end{umlpackage}
    \end{tikzpicture}
    \caption{UML Diagram for \texttt{AInteger}}
    \label{fig:uml-ainteger}
\end{figure}

% Creating UML diagram for AFloat
\begin{figure}[h]
    \centering
    \begin{tikzpicture}
        \begin{umlpackage}{arbitraryarithmetic}
            \umlclass{AFloat}{
                - integerPart: String \\
                - fractionalPart: String \\
                - isNegative: boolean \\
                - MAX\_PRECISION: int = 30 \\
                \\
                + AFloat() \\
                + AFloat(s: String) \\
                + AFloat(other: AFloat) \\
                + parse(s: String): AFloat \\
                + toString(): String \\
                + add(b: AFloat): AFloat \\
                + subtract(b: AFloat): AFloat \\
                + multiply(b: AFloat): AFloat \\
                + divide(b: AFloat): AFloat \\
                - removeLeadingZeros(s: String): String \\
                - removeTrailingZeros(s: String): String
            }{}
            \umlclass[below of=AFloat, yshift=-1cm]{AInteger}{...}{}
            \umluniaggreg[geometry=|-|, anchor1=180, anchor2=0]{AFloat}{AInteger}
        \end{umlpackage}
    \end{tikzpicture}
    \caption{UML Diagram for \texttt{AFloat} (depends on \texttt{AInteger})}
    \label{fig:uml-afloat}
\end{figure}

% Creating UML diagram for MyInfArith
\begin{figure}[h]
    \centering
    \begin{tikzpicture}
        \umlclass{MyInfArith}{
            \\
            + main(args: String[]): void
        }{}
        \begin{umlpackage}[right of=MyInfArith, xshift=4cm]{arbitraryarithmetic}
            \umlclass{AInteger}{...}{}
            \umlclass[below of=AInteger]{AFloat}{...}{}
        \end{umlpackage}
        \umluniaggreg[geometry=|-|, anchor1=0, anchor2=180]{MyInfArith}{AInteger}
        \umluniaggreg[geometry=|-|, anchor1=0, anchor2=180]{MyInfArith}{AFloat}
    \end{tikzpicture}
    \caption{UML Diagram for \texttt{MyInfArith} (uses \texttt{AInteger} and \texttt{AFloat})}
    \label{fig:uml-myinfarith}
\end{figure}

% Introducing the README section
\section{README: How to Use the Library}
This section provides instructions for users to build and run the library, adapted from our \texttt{README.md}.

\subsection{Prerequisites}
\begin{itemize}
    \item Java Development Kit (JDK) 17 or higher.
    \item Python 3 for running the automation script.
    \item Apache Ant for building the project.
    \item Docker (optional) for containerized execution.
\end{itemize}

\subsection{Project Structure}
\begin{lstlisting}[language=bash]
project/
- arbitraryarithmetic/
  - AInteger.java
  - AFloat.java
- MyInfArith.java
- compile_and_run.py
- build.xml
- Dockerfile
- README.md
- report.tex
\end{lstlisting}

\subsection{Building the Project}
\begin{enumerate}
    \item \textbf{Using Ant}:
    \begin{lstlisting}[language=bash]
ant
    \end{lstlisting}
    This compiles the Java files and creates \texttt{build/aarithmetic.jar} and \texttt{build/classes/} with compiled classes.

    \item \textbf{Using Docker}:
    Build the Docker image:
    \begin{lstlisting}[language=bash]
docker build -t arithmetic-app .
    \end{lstlisting}
\end{enumerate}

\subsection{Running the Program}
\begin{enumerate}
    \item \textbf{Using the Python Script (Local)}:
    \begin{lstlisting}[language=bash]
python3 compile_and_run.py int add 123 456
    \end{lstlisting}
    This compiles and runs \texttt{MyInfArith} with arguments \texttt{int add 123 456}.

    \item \textbf{Directly Running Java (Local)}:
    \begin{lstlisting}[language=bash]
java -cp build/aarithmetic.jar MyInfArith int add 123 456
    \end{lstlisting}

    \item \textbf{Using Docker}:
    \begin{lstlisting}[language=bash]
docker run --rm arithmetic-app int add 123 456
    \end{lstlisting}
\end{enumerate}

\subsection{Example Outputs}
\begin{itemize}
    \item \texttt{python3 compile\_and\_run.py int add 123 456}:
    \begin{lstlisting}
123 add 456 = 579
    \end{lstlisting}
    \item \texttt{python3 compile\_and\_run.py float mul 123.456 -78.123}:
    \begin{lstlisting}
123.456 mul -78.123 = -9643.237576288
    \end{lstlisting}
\end{itemize}

% Introducing the Git commits section
\section{Snapshot of Git Commits}
Below is a snapshot of our Git commit history, showing the project’s development stages:
\begin{lstlisting}
5d3f9a2 Stage 5: Final submission with report and README
4e2b8c1 Stage 4: Added Ant build.xml and Dockerfile
3a1d7f0 Stage 3: Added MyInfArith.java and compile_and_run.py
2b6c5e9 Stage 2: Implemented AFloat.java
1c4a3d8 Stage 1: Implemented AInteger.java
0f9e2c7 Initial commit: Project setup
\end{lstlisting}
Each stage was tagged (e.g., \texttt{v1.0} to \texttt{v5.0}) for clarity.

% Introducing the limitations section
\section{Limitations of the Library}
\begin{itemize}
    \item \textbf{Performance}: Division in \texttt{AInteger} uses repeated subtraction, which is slow for large numbers. A more efficient algorithm (e.g., long division) could improve speed.
    \item \textbf{Decimal Precision}: \texttt{AFloat} truncates to 30 decimal digits, which may not suffice for applications needing higher precision.
    \item \textbf{Error Handling}: Errors (e.g., invalid input, division by zero) are handled with console messages rather than exceptions, limiting robustness.
    \item \textbf{No Advanced Operations}: The library supports basic arithmetic (+, -, *, /) but lacks operations like exponentiation or square roots.
\end{itemize}

% Introducing the verification approach
\section{Verification Approach}
We verified the library through the following methods:
\begin{itemize}
    \item \textbf{Unit Testing}: Manually tested edge cases in \texttt{Main.java} and \texttt{compile\_and\_run.py}:
    \begin{itemize}
        \item Zero cases: \texttt{0 + 0}, \texttt{0 / 5}.
        \item Negative numbers: \texttt{-123 + 456}, \texttt{-5 * -3}.
        \item Large numbers: \texttt{123456789012345 + 987654321}.
        \item Invalid inputs: \texttt{"abc"}, \texttt{"12.34.56"} (handled by \texttt{AInteger}/\texttt{AFloat}).
    \end{itemize}
    \item \textbf{Mental Simulation}: Walked through arithmetic algorithms (e.g., digit-by-digit addition) to ensure correctness.
    \item \textbf{Automation Testing}: Used \texttt{compile\_and\_run.py} to run multiple test cases, verifying outputs against expected results.
    \item \textbf{Docker Testing}: Ran the Docker container to ensure the library works in a containerized environment.
\end{itemize}
All operations preserve precision, with \texttt{AFloat} truncating to 30 decimal digits as specified.

% Introducing the key learnings section
\section{Key Learnings}
\begin{itemize}
    \item \textbf{Arbitrary-Precision Arithmetic}: Learned to implement string-based arithmetic to overcome Java’s numeric limits, handling carries and decimal points manually.
    \item \textbf{Object-Oriented Design}: Understood the importance of separating concerns (\texttt{AInteger} vs. \texttt{AFloat}) and reusing code (\texttt{AFloat} uses \texttt{AInteger}).
    \item \textbf{Tools and Automation}: Gained experience with Ant for building, Docker for containerization, and Python for scripting, streamlining development and deployment.
    \item \textbf{Documentation}: Realized the value of clear documentation at all levels (code, report, README) for usability and maintainability.
    \item \textbf{Testing}: Learned the importance of testing edge cases (e.g., division by zero, large numbers) to ensure robustness.
\end{itemize}

\end{document}